\scanpage{77}
% archive foliation: 76

Alors, définissant $T_0$ par (8) (pour $t\in\Sigma_J^*$ donné) on a un isom canonique
\[
(12)\qquad T_0^{\otimes2}\simeq T_{\Sigma,t}
\]
Ceci posé, on appellera \uline{épinglage de Legendre}
\deleted{\unsure{car il faut se rappeler que}}
\deleted{$\unsure{T_0}$ \unsure{ne correspond}}
de $E$, subordonné au carré $Q_0$ (muni des données (11)) \deleted{tout} \inserted{un} isomorphisme
\[
(13)\qquad T=T_{E,0}\xleftarrow{\substack{\kappa\\\sim}}T_0
\]
compatible avec les \unsure{hom. carrure}, i.e. rendant commutatif le diagramme
\diagram{$(14)\quad\begin{tikzcd}[ampersand replacement=\&,column sep=small,row sep=large]
T_0^{\otimes2}\arrow[rr,"{\substack{\kappa^{\otimes2}\\\sim}}"]\arrow[dr,"(12)"]\arrow[d,dashed,no head]\&\&T^{\otimes2}\arrow[dl,"(5)"]\\
{}\&T_{\Sigma,t}\&
\end{tikzcd}$}

Il y en a exactement deux, différant par le signe i.e. formant un torseur sous $\mu_2$.

Bien entendu, on veut pouvoir définir la notion d'épinglage de $E$ (en termes d'un $Q_0$ donné une fois pour toutes) pour $E$ variable --- mais alors $\underline\omega=\underline\omega_J$ n'est pas encore déterminé (puisque $J$ ne l'est pas). Donc la donnée d'un épinglage de Legendre sera définie comme \unsure{le couple|les couples} \inserted{$(\alpha,\kappa)$} formé d'un isom.
\[
(15)\qquad\alpha:\underline\omega_J\xleftarrow{\sim}D_{Q_0}=Q_0/\{\pm1\}
\]
\uline{et} d'un isom. $\kappa$ (13) satisfaisant la compatibilité (14).
